\documentclass[11pt]{article}
\usepackage{amsmath,amssymb,amsthm,hyperref}
\newtheorem{theorem}{Theorem}
\newtheorem{lemma}[theorem]{Lemma}
\newcommand{\E}{\mathbb E}
\newcommand{\Var}{\operatorname{Var}}
\title{The logarithmic factor in the large-block Gaussian limit}
\author{Research note; independently reviewed by two other agents}
\date{30 September 2026}
\begin{document}
\maketitle

Let $R^{(n,m)}_\sigma=n!P(\sigma)$ be the likelihood ratio of a target
order for $m$ independent uniformly relabelled palindrome blocks.
This note treats an iterated limit: first $m\to\infty$ for each fixed $n$,
then $n\to\infty$. It makes no claim that the Gaussian approximation is
uniform at the threshold $m\asymp n^{13/10}(\log n)^{1/5}$.

\begin{theorem}\label{thm:main}
For every fixed $n\geq3$, the finite vector
\[
 \bigl(m^{5/2}(R^{(n,m)}_\sigma-1):\sigma\in S_n\bigr)
\]
converges in distribution as $m\to\infty$ to a centered Gaussian vector
$(G_\sigma:\sigma\in S_n)$. The expectations of the maximum absolute
coordinates also converge. Moreover, for absolute constants $c,C>0$,
\[
 c n^{13/4}\sqrt{\log n}\leq
 \lim_{m\to\infty}m^{5/2}\E\max_{\sigma\in S_n}|R^{(n,m)}_\sigma-1|
 \leq C n^{13/4}\sqrt{\log n}
\]
for sufficiently large $n$.
\end{theorem}

\section{The Gaussian limit}
Put $d=n-3$, $B_n=\binom n3$, and
\[
 w_a(p)=\binom da p^a(1-p)^{d-a},\qquad
 \delta(l,x,r)=\frac12-\frac32\mathbf1\{x<\min(l,r)\}.
\]
Let $Y=(Y_1,\ldots,Y_n)$ consist of independent uniform $[0,1]$ variables,
and define
\[
 f_\sigma(p,Y)=\sum_{a=0}^{d}w_a(p)
       \delta(Y_{\sigma(a)},Y_{\sigma(a+1)},Y_{\sigma(a+2)}),
\]
where target positions are numbered from zero. Each $f_\sigma$ is centered
as a function of $Y$. The triple projection is exactly
\[
 S_\sigma=B_nm^{-3}(1-1/m)^d
             \sum_{i=1}^m f_\sigma((i-1)/(m-1),Y_i),
\]
with independent $Y_i$. Its scaled summands satisfy the Lindeberg condition
for fixed $n$, since they are uniformly $O_n(m^{-1/2})$.
Their covariance sums converge by ordinary Riemann integration. Therefore
the multivariate Lindeberg central limit theorem gives the centered
Gaussian limit with covariance
\begin{equation}\label{eq:cov}
 \E G_\sigma G_\tau
 =B_n^2\int_0^1\E[f_\sigma(p,Y)f_\tau(p,Y)]\,dp.
\end{equation}

For clarity, the other terms vanish after scaling. At Hoeffding level
$\ell$, the blockwise centered kernel on any specified set of $\ell$ blocks
is $O_n(m^{-3\ell})$. Symmetrization, or directly orthogonality of distinct
canonical kernels for $L^2$, bounds the full level by
$O_n(m^{-5\ell/2})$ in $L^2$. Levels $\ell\geq2$ are therefore negligible.
In level one the segments of length at least four give
$O_n(m^{-7/2})$ in $L^2$. Thus
\[
 m^{5/2}(R^{(n,m)}_\sigma-1-S_\sigma)\longrightarrow0
 \quad\hbox{in }L^2.
\]
There are finitely many targets for fixed $n$, and their scaled coordinates
have bounded second moments. Their maximum is consequently uniformly
integrable, proving convergence of the expected maximum absolute value.

Equivalently, let $W$ be an isonormal Gaussian process over the Hilbert
space $L^2([0,1]\times[0,1]^n,dp\,dY)$. Then
\[
 G_\sigma=B_n W(f_\sigma).
\]
Only a finite-dimensional subspace is needed here, so this is simply a
convenient representation of a finite Gaussian vector with covariance
\eqref{eq:cov}.

\section{The upper bound}
The centered valley sequence in a fixed target has covariances
$1/2,-1/4,1/20$ at distances zero, one, and two, and zero beyond two.
By Cauchy--Schwarz applied to shifted weight sequences,
\[
 \Var\left(\sum_aw_a(p)\delta_a\right)
 \leq\frac{11}{10}\sum_aw_a(p)^2.
\]
The exact beta-integral identity
\[
 \int_0^1\sum_aw_a(p)^2\,dp
 =\frac{4^d}{(2d+1)\binom{2d}d}\asymp n^{-1/2}
\]
therefore gives $\Var G_\sigma\leq Cn^{11/2}$.
The elementary Gaussian exponential-moment bound for a maximum over $n!$
coordinates now implies
\[
 \E\max_\sigma|G_\sigma|
 \leq Cn^{11/4}\sqrt{\log(2n!)}
 \leq Cn^{13/4}\sqrt{\log n}.
\]
No independence among these coordinates is required.

\section{A canonical three-variable component}
Let $L,X,R$ be independent uniform variables, and let $u(L,X,R)$ be the
component of $\delta(L,X,R)$ centered in each of its three variables:
\[
 u=\delta-\E(\delta\mid L,X)-\E(\delta\mid X,R)
            -\E(\delta\mid L,R)
       +\E(\delta\mid L)+\E(\delta\mid X)+\E(\delta\mid R).
\]
It has zero conditional expectation given any two variables.
Its variance is
\begin{equation}\label{eq:kappa}
 \E u^2=\frac1{20}.
\end{equation}
Here is an explicit check. With
$g(t)=\tfrac34(1-t)^2-\tfrac14$, the one-variable projections are
$g(L),-2g(X),g(R)$ and $\E g^2=1/20$.
The variances of the two-variable conditional expectations are
\[
 \Var\E(\delta\mid L,R)=\frac18,
 \qquad
 \Var\E(\delta\mid L,X)=\Var\E(\delta\mid X,R)=\frac5{16}.
\]
Orthogonality of the Hoeffding components gives
$\E u^2=\tfrac12+6/20-\tfrac18-2(5/16)=1/20$.

\section{An independent Gaussian assignment matrix inside the process}
Choose $N\asymp n$ movable positions in the middle third of a base target,
spaced by at least three positions. Keep all other labels fixed as anchors,
and allow the $N$ movable labels to permute arbitrarily among the movable
positions. Let $\mathcal F$ be this family of $N!$ target orders.
Each movable position $v$ has two fixed adjacent anchor labels
$\lambda_v,\rho_v$. These anchor pairs are disjoint between positions,
and no anchor is a movable label.

For a movable label $b$, define the Gaussian entry
\[
 Z_{v,b}=B_nW\bigl(w_{v-1}(p)
                     u(Y_{\lambda_v},Y_b,Y_{\rho_v})\bigr).
\]
These $N^2$ entries are mutually independent. Indeed, distinct entries use
different triples of underlying variables, with at most two variables in
common. Integrating in a variable present in only one triple makes their
inner product zero by the canonical property of $u$. Joint Gaussianity
turns this orthogonality into independence.
For each row $v$, all entries have the same variance, and
\[
 \Var Z_{v,b}=
 \frac{B_n^2}{20}\int_0^1w_{v-1}(p)^2\,dp\asymp n^{9/2}.
\]
The constants are uniform over the chosen central positions: the integral
is $\binom d{v-1}^2/((2d+1)\binom{2d}{2v-2})\asymp n^{-3/2}$.

If a target $\sigma\in\mathcal F$ assigns label $\pi(v)$ to position $v$,
its orthogonal projection onto the span of these Gaussian entries is
precisely $\sum_vZ_{v,\pi(v)}$. To check this, a target triple can have a
nonzero inner product with the defining canonical kernel only if it
contains all three labels $\lambda_v,b,\rho_v$. The two fixed anchors
are separated by position $v$, so the only such consecutive triple has
center label $b$ assigned to $v$. Its coefficient is exactly $w_{v-1}$.
Consequently,
\[
 G_\sigma=\sum_vZ_{v,\pi(v)}+H_\sigma,
\]
where the entire centered Gaussian residual vector $(H_\sigma)$ is
independent of the matrix $Z$.

Conditional Jensen for the maximum gives
\[
 \E\max_{\sigma\in\mathcal F}G_\sigma
 \geq\E\max_{\pi\in S_N}\sum_vZ_{v,\pi(v)}.
\]
A greedy assignment proves a lower bound for the right side: process rows
one by one, and select in each row the largest entry among the remaining
columns. The remaining set depends only on previous rows, so it is
independent of the current row. For the first $N/2$ rows, at least $N/2$
independent centered Gaussians of common standard deviation
$\asymp n^{9/4}$ remain. Their expected maximum is at least
$c n^{9/4}\sqrt{\log N}$. Later row increments have nonnegative expectation.
Thus
\[
 \E\max_{\sigma\in S_n}|G_\sigma|
 \geq cNn^{9/4}\sqrt{\log N}
 \geq c n^{13/4}\sqrt{\log n}.
\]
Together with the Gaussian upper bound and the convergence of expectations,
this proves Theorem~\ref{thm:main}.

\paragraph{Interpretation.}
The logarithmic factor in the uniform upper estimate reflects real
behavior of this random construction in its Gaussian limit. Extending this
lower bound uniformly down to $m\asymp n^{13/10}(\log n)^{1/5}$ would require
additional quantitative comparison or direct assignment estimates.
The separate adjacent-swap argument already proves a uniform lower bound
without the logarithm in the regime $m\geq Cn^{13/10}$.
\end{document}
